Este anexo consolida la documentaci\'on metodol\'ogica y de c\'odigo de los tres
notebooks principales de la entrega: Markowitz vs Benchmark, Black-Litterman y
Simulaci\'on Monte Carlo. La reorganizaci\'on mantiene la informaci\'on sustantiva
de los anexos metodol\'ogicos originales, pero la presenta en un solo anexo y con
la misma estructura para cada t\'ecnica: resumen, fundamentos metodol\'ogicos,
supuestos del experimento, estructura del notebook, KPIs utilizados, resultados
del experimento pandemia, rebalanceo, splits aleatorios, evaluaci\'on robusta e
implicancias y conclusiones.

% ============================================================
\section*{Markowitz vs Benchmark}
% ============================================================

\subsection*{Resumen}

Este apartado resume el funcionamiento, los fundamentos y los principales
resultados del notebook
\texttt{markowitz\_vs\_benchmark\_rebalanceo\_pandemia.ipynb}. El objetivo del
experimento fue comparar portafolios construidos mediante optimizaci\'on
media-varianza de Markowitz contra benchmarks simples y variantes te\'oricas,
evaluando no solo el retorno obtenido, sino tambi\'en la robustez del modelo ante
cambios en tres decisiones metodol\'ogicas: horizonte de calibraci\'on y
evaluaci\'on futura, inclusi\'on o exclusi\'on del periodo de pandemia 2020--2022
en los datos de entrenamiento, y uso de rebalanceo semestral. Adem\'as, se
incorpora como referencia un benchmark top-20 equiponderado por capitalizaci\'on
burs\'atil.

El caso principal utiliza 601 activos, 7 a\~nos hist\'oricos de calibraci\'on,
2 a\~nos futuros fuera de muestra, tasa libre de riesgo anual de 2\%, capital
inicial de 1\,000, cota m\'axima de peso de 20\% para la frontera eficiente y
rebalanceo cada 126 d\'ias h\'abiles, equivalente aproximadamente a medio a\~no
burs\'atil. La evaluaci\'on es fuera de muestra: los pesos se estiman con datos
hist\'oricos y luego se aplican al periodo futuro. Esto permite distinguir entre
desempe\~no te\'orico dentro de la muestra y desempe\~no real en datos no usados
para optimizar.

\subsection*{Fundamentos metodol\'ogicos}

El modelo de Markowitz representa cada portafolio mediante un vector de pesos
$w$, un vector de retornos esperados $\mu$ y una matriz de covarianzas
$\Sigma$. En el notebook, $\mu$ se estima como promedio diario hist\'orico y
$\Sigma$ como covarianza diaria de retornos. Para cada perfil de riesgo se
resuelve el problema:

\begin{equation}
\max_w \; \mu^\top w - \frac{1}{2}\gamma w^\top \Sigma w,
\end{equation}

sujeto a restricciones de presupuesto, no negatividad, peso m\'aximo y, cuando
corresponde, una cota de volatilidad:

\begin{equation}
\sum_i w_i = 1, \qquad 0 \leq w_i \leq w_{\max}, \qquad
w^\top \Sigma w \leq \sigma^2_{\max}.
\end{equation}

La restricci\'on $w_i \geq 0$ impide ventas cortas, la suma igual a uno invierte
todo el capital y el l\'imite por activo evita concentraci\'on extrema. El
par\'ametro $\gamma$ controla la aversi\'on al riesgo: valores altos penalizan
m\'as la varianza y producen carteras conservadoras, mientras que valores bajos
permiten mayor exposici\'on a retorno esperado.

La frontera eficiente se obtiene resolviendo, para distintos retornos objetivo,
el portafolio de m\'inima varianza que alcanza al menos dicho retorno:

\begin{equation}
\min_w \; w^\top \Sigma w.
\end{equation}

El portafolio denominado \textit{Ideal de mercado} se selecciona como el punto
de la frontera eficiente con mayor ratio de Sharpe,

\begin{equation}
Sharpe = \frac{R_p - R_f}{\sigma_p}.
\end{equation}

Esta es una aproximaci\'on discreta al portafolio de tangencia, porque el c\'odigo
no maximiza directamente el cociente de Sharpe, sino que elige el mejor punto
entre 30 puntos calculados sobre la frontera.

El benchmark base selecciona los 20 primeros activos del universo ordenado por
\texttt{marketCap} y asigna pesos iguales. Esta regla no estima $\mu$ ni
$\Sigma$, por lo que es menos sensible al error de par\'ametros. Adem\'as, se
define \texttt{b\_teorico}, que toma la volatilidad hist\'orica del benchmark y
construye un portafolio Markowitz que maximiza retorno esperado sujeto a esa
misma volatilidad. Su objetivo es responder qu\'e portafolio te\'orico podr\'ia
obtenerse si se acepta el mismo nivel de riesgo que el benchmark.

Para estabilizar la matriz de covarianzas, el notebook aplica \textit{shrinkage}
de 20\%, mezclando 80\% de la covarianza completa con 20\% de su diagonal:

\begin{equation}
\Sigma_{\text{shrink}} = (1-\lambda)\Sigma + \lambda \, \text{diag}(\Sigma).
\end{equation}

Luego se utiliza \texttt{nearest\_psd} para forzar que la matriz sea
semidefinida positiva, requisito t\'ecnico para resolver problemas cuadr\'aticos
convexos con CVXPY. El rebalanceo divide el periodo futuro en segmentos de
126 d\'ias h\'abiles; al inicio de cada segmento se reestima el portafolio usando
la historia inicial y los retornos futuros ya observados hasta ese momento.
As\'i, el modelo incorpora informaci\'on nueva sin usar datos futuros no
disponibles al momento de la decisi\'on.

\subsection*{Supuestos del experimento}

El experimento supone que los precios de cierre y dividendos de los CSV
representan adecuadamente la informaci\'on disponible para cada activo. Los
retornos diarios se calculan como retornos totales, incorporando dividendos, por
lo que se considera que estos forman parte del rendimiento econ\'omico del
inversionista. Al construir la matriz de retornos se utiliza
\texttt{join="inner"}, conservando solo fechas comunes a todos los activos
seleccionados; esto facilita la estimaci\'on matricial, aunque asume que eliminar
fechas no comunes no introduce sesgos relevantes.

Desde el punto de vista de mercado, las carteras son \textit{long-only},
invierten el 100\% del capital disponible, no usan apalancamiento y permiten
divisibilidad perfecta de activos. El modelo omite costos de transacci\'on,
impuestos, \textit{spreads} bid-ask, restricciones de liquidez, impacto de
mercado y \textit{slippage}. Por ello, los resultados deben interpretarse como
una aproximaci\'on te\'orica del desempe\~no y no como una simulaci\'on completa de
ejecuci\'on real. La tasa libre de riesgo se mantiene constante en 2\% anual.

En t\'erminos estad\'isticos, se asume que los retornos hist\'oricos contienen
informaci\'on \'util sobre retornos y riesgos futuros, aunque los propios
resultados muestran que esta estabilidad es imperfecta. La anualizaci\'on usa
252 d\'ias h\'abiles; el retorno se anualiza por capitalizaci\'on compuesta y la
volatilidad mediante ra\'iz del tiempo. Si una cota de volatilidad vuelve
infactible un perfil, el c\'odigo relaja esa cota manteniendo las restricciones
restantes. Finalmente, en el experimento pandemia el periodo futuro fuera de
muestra se mantiene fijo entre escenarios, por lo que las diferencias observadas
se atribuyen a la informaci\'on usada para calibrar y no a cambios en el mercado
futuro evaluado.

\subsection*{Estructura del notebook}

El notebook se organiza en bloques secuenciales. Primero define la
documentaci\'on inicial, requisitos, estructura de carpetas y uso de archivos
cacheados mediante \texttt{FORCE\_RECOMPUTE}. Luego importa librer\'ias, fija
rutas, define \texttt{TRADING\_DAYS\_PER\_YEAR = 252},
\texttt{REBALANCE\_DAYS = 126}, delimita el periodo pand\'emico y crea los
escenarios \texttt{sin\_pandemia} y \texttt{con\_pandemia}. Posteriormente carga
el script base \texttt{markowitz\_cvxpy\_benchmark\_teorico.py} como m\'odulo
local, reutilizando funciones de lectura, selecci\'on de universo, estimaci\'on
de par\'ametros, optimizaci\'on, frontera eficiente, benchmark y evaluaci\'on.

El experimento principal fija \texttt{DEFAULT\_BENCHMARK\_SIZE = 20},\\
\texttt{DEFAULT\_MARKET\_UNIVERSE\_SIZE = 601},
\texttt{DEFAULT\_FRONTIER\_POINTS = 30},
\texttt{DEFAULT\_FRONTIER\_MAX\_WEIGHT = 0.20},
\texttt{DEFAULT\_RISK\_FREE\_RATE = 0.02},
\texttt{DEFAULT\_INITIAL\_CAPITAL = 1000.0},
\texttt{DEFAULT\_SHRINKAGE = 0.20},\\
\texttt{HISTORICAL\_YEARS = 7} y \texttt{FUTURE\_YEARS = 2}. La configuraci\'on
\texttt{h7\_f2} se usa como caso central porque la validaci\'on robusta posterior
la identifica como el horizonte con mejor desempe\~no agregado.

Las funciones auxiliares del notebook permiten cargar datasets por escenario,
calcular retornos totales, generar ventanas de rebalanceo, definir la vista de
entrenamiento disponible en cada segmento, resolver portafolios por instante y
evaluar el periodo futuro completo. La clase \texttt{ScenarioDataset} agrupa la
informaci\'on de cada escenario, mientras que
\texttt{evaluate\_rebalanced\_window} coordina el experimento: crea el split
hist\'orico/futuro, divide el futuro en segmentos, reestima pesos en cada
rebalanceo, calcula retornos realizados y guarda m\'etricas globales, por
segmento y por punto de frontera.

La validaci\'on inicial verifica rutas, carga de datos y funcionamiento de
ventanas. En una prueba reducida se obtienen 2\,717 observaciones para el
escenario sin pandemia y 3\,473 para el escenario con pandemia, diferencia
esperable porque la segunda base conserva 2020--2022. En el caso principal se
cargan o generan cuatro CSV:\\
\texttt{rebalance\_summary\_h7\_f2\_601activos.csv},\\
\texttt{rebalance\_segments\_h7\_f2\_601activos.csv},\\
\texttt{segment\_summary\_h7\_f2\_601activos.csv} y\\
\texttt{frontier\_segments\_h7\_f2\_601activos.csv}.

\begin{table}[H]
\centering
\caption{Dimensi\'on de resultados cacheados del experimento principal.}
\label{tab:anexo4_cache}
\small
\begin{tabular}{@{}lrrr@{}}
\toprule
Escenario & Portafolios & Filas segmento & Puntos frontera \\
\midrule
Sin pandemia & 9 & 36 & 119 \\
Con pandemia & 9 & 36 & 119 \\
\bottomrule
\end{tabular}
\end{table}

Los 9 portafolios corresponden a cinco perfiles de riesgo, Benchmark,
\texttt{b\_teorico}, M\'inima varianza e Ideal de mercado. Las 36 filas por
escenario equivalen a 9 portafolios evaluados en 4 segmentos semestrales. Al
exportar resultados, el notebook genera 18 filas de resumen global, 72 filas de
segmentos rebalanceados, 72 filas de resumen por segmento y 238 puntos de
frontera.

\subsection*{KPIs utilizados}

El desempe\~no se eval\'ua mediante indicadores rebalanceados fuera de muestra.
El retorno total rebalanceado mide la ganancia acumulada durante todo el periodo
futuro:

\begin{equation}
R_{\text{total}} = \prod_t (1+r_t)-1.
\end{equation}

El retorno anualizado para los dos a\~nos futuros se calcula como:

\begin{equation}
R_{\text{anual}} = (1+R_{\text{total}})^{1/2}-1.
\end{equation}

La volatilidad diaria corresponde a la desviaci\'on est\'andar de los retornos
concatenados de todos los segmentos y se anualiza mediante:

\begin{equation}
\sigma_{\text{anual}} = \sigma_{\text{diaria}}\sqrt{252}.
\end{equation}

El Sharpe rebalanceado mide retorno excedente por unidad de riesgo:

\begin{equation}
Sharpe = \frac{R_{\text{anual}} - 0.02}{\sigma_{\text{anual}}}.
\end{equation}

Tambi\'en se calculan m\'aximo \textit{drawdown}, CVaR diario al 95\%, n\'umero
de activos con peso positivo, mayor peso individual, cinco principales
posiciones, riqueza final futura desde capital inicial 1\,000, Sharpe te\'orico
y Sharpe real futuro. Esta \'ultima comparaci\'on es central para diagnosticar
\textit{overfitting}, ya que contrasta la eficiencia esperada durante el
entrenamiento con la eficiencia realizada fuera de muestra.

\subsection*{Resultados del experimento pandemia}

La comparaci\'on entre escenarios calcula retorno total, retorno anual,
volatilidad anual, Sharpe y diferencias
\texttt{con\_pandemia - sin\_pandemia}. El resultado central se resume en la
Tabla~\ref{tab:anexo4_comparacion_pandemia}.

\begin{table}[H]
\centering
\caption{Comparaci\'on de portafolios con y sin pandemia en entrenamiento.}
\label{tab:anexo4_comparacion_pandemia}
\scriptsize
\begin{adjustbox}{max width=\textwidth}
\begin{tabular}{@{}lrrrrrr@{}}
\toprule
Portafolio & Ret. sin pand. & Sharpe sin pand. & Ret. con pand. & Sharpe con pand. & Delta ret. & Delta Sharpe \\
\midrule
Muy conservador & 58,0\% & 2,208 & 54,5\% & 2,023 & $-$3,4\% & $-$0,184 \\
Conservador & 67,8\% & 2,425 & 55,5\% & 1,972 & $-$12,3\% & $-$0,453 \\
Neutro & 82,9\% & 2,413 & 60,5\% & 1,756 & $-$22,5\% & $-$0,657 \\
Arriesgado & 91,6\% & 1,733 & 58,4\% & 1,167 & $-$33,2\% & $-$0,566 \\
Muy arriesgado & 102,7\% & 1,412 & 44,0\% & 0,549 & $-$58,7\% & $-$0,863 \\
Benchmark & 101,4\% & 2,052 & 101,4\% & 2,052 & $-$0,0\% & $-$0,000 \\
\texttt{b\_teorico} & 104,5\% & 1,515 & 56,7\% & 0,966 & $-$47,8\% & $-$0,548 \\
M\'inima varianza & 56,6\% & 2,110 & 54,6\% & 2,016 & $-$2,0\% & $-$0,094 \\
Ideal de mercado & 76,5\% & 2,507 & 59,5\% & 1,399 & $-$16,9\% & $-$1,108 \\
\bottomrule
\end{tabular}
\end{adjustbox}
\end{table}

Sin pandemia, el mejor portafolio es Ideal de mercado, con Sharpe 2,507 y
retorno total de 76,5\%. Con pandemia, el mejor resultado corresponde al
Benchmark, con Sharpe 2,052 y retorno de 101,4\%. Incluir el periodo
2020--2022 deteriora casi todos los portafolios optimizados, especialmente
aquellos m\'as agresivos o dependientes de estimaciones de retorno esperado. El
Benchmark permanece invariante porque su regla no utiliza $\mu$ ni $\Sigma$,
sino una selecci\'on top-20 equiponderada.

\begin{table}[H]
\centering
\caption{Mayores ca\'idas de Sharpe al incluir pandemia en entrenamiento.}
\label{tab:anexo4_delta_sharpe}
\small
\begin{tabular}{@{}lr@{}}
\toprule
Portafolio & Delta Sharpe \\
\midrule
Ideal de mercado & $-$1,108 \\
Muy arriesgado & $-$0,863 \\
Neutro & $-$0,657 \\
Arriesgado & $-$0,566 \\
\texttt{b\_teorico} & $-$0,548 \\
\bottomrule
\end{tabular}
\end{table}

La interpretaci\'on es que los datos de pandemia alteran las estimaciones de
retornos y covarianzas. El efecto es m\'as fuerte en carteras agresivas o de
tangencia, que dependen de diferencias peque\~nas en $\mu$. En cambio, M\'inima
varianza es la cartera Markowitz m\'as estable, con una ca\'ida de Sharpe de solo
2,110 a 2,016, porque depende principalmente de covarianzas y no de
predicciones de retorno individual.

\subsection*{Rebalanceo}

El rebalanceo semestral mejora el Sharpe en Conservador, Ideal de mercado,
M\'inima varianza y Neutro, pero deteriora Muy arriesgado, Arriesgado,
\texttt{b\_teorico} y Benchmark frente a una versi\'on base sin rebalanceo. Por
lo tanto, el rebalanceo funciona como un mecanismo de actualizaci\'on y control,
pero no garantiza mejor desempe\~no. Puede reducir deriva de pesos e incorporar
informaci\'on reciente, aunque tambi\'en puede vender activos ganadores demasiado
pronto o amplificar ruido en estimaciones nuevas.

\subsection*{Splits aleatorios}

El experimento de splits aleatorios no recalcula Markowitz, sino que lee
salidas preexistentes de
\texttt{outputs\_markowitz\_cvxpy\_splits\_aleatorios}. El dise\~no considera
una ventana base de 9 a\~nos, 28 combinaciones train/test con suma menor o igual
a 9, 70 ventanas aleatorias con semilla 42, 560 evaluaciones de portafolio,
601 activos y 1\,392 puntos de frontera.

\begin{table}[H]
\centering
\caption{Resumen del experimento de splits aleatorios.}
\label{tab:anexo4_splits_resumen}
\small
\begin{tabular}{@{}lr@{}}
\toprule
Indicador & Resultado \\
\midrule
Combinaciones train/test & 28 \\
Ventanas aleatorias & 70 \\
Evaluaciones de portafolios & 560 \\
Activos & 601 \\
Puntos de frontera & 1\,392 \\
\bottomrule
\end{tabular}
\end{table}

El mejor split promedio global es \texttt{h5\_f2}, con Sharpe futuro promedio de
2,34. Para el Benchmark, el mejor horizonte particular es \texttt{h6\_f2}, con
Sharpe futuro medio de 2,615. Esta sensibilidad muestra que las conclusiones no
deben depender de un \'unico split.

\begin{table}[H]
\centering
\caption{Promedio y variabilidad del Sharpe en splits aleatorios.}
\label{tab:anexo4_splits_sharpe}
\small
\begin{tabular}{@{}lrr@{}}
\toprule
Portafolio & Sharpe medio & Desv. est. Sharpe \\
\midrule
Benchmark & 2,057 & 0,381 \\
M\'inima varianza & 1,940 & 0,250 \\
Muy conservador & 1,885 & 0,288 \\
Conservador & 1,799 & 0,334 \\
Ideal de mercado & 1,799 & 0,421 \\
Neutro & 1,781 & 0,413 \\
Arriesgado & 1,644 & 0,449 \\
Muy arriesgado & 1,625 & 0,458 \\
\bottomrule
\end{tabular}
\end{table}

En promedio, el Benchmark supera a los portafolios Markowitz. M\'inima varianza
aparece como la versi\'on Markowitz m\'as robusta, ya que combina buen Sharpe con
baja variabilidad.

\subsection*{Evaluaci\'on robusta}

La validaci\'on robusta con ventanas m\'oviles eval\'ua 6 horizontes y 18 ventanas
fuera de muestra; sus resultados respaldan \texttt{h7\_f2} como caso principal.

\begin{table}[H]
\centering
\caption{Sharpe futuro medio en validaci\'on robusta por ventanas m\'oviles.}
\label{tab:anexo4_validacion_robusta}
\small
\begin{tabular}{@{}lr@{}}
\toprule
Horizonte & Sharpe futuro medio \\
\midrule
\texttt{h7\_f2} & 2,32 \\
\texttt{h6\_f3} & 2,07 \\
\texttt{h5\_f4} & 1,88 \\
\texttt{h4\_f5} & 1,66 \\
\texttt{h3\_f6} & 1,60 \\
\texttt{h2\_f7} & 1,35 \\
\bottomrule
\end{tabular}
\end{table}

El resultado indica que 7 a\~nos de historia y 2 a\~nos de futuro equilibran
estabilidad de par\'ametros y relevancia temporal. Adem\'as, en los splits
aleatorios el Ideal de mercado presenta Sharpe te\'orico medio de 5,01 y Sharpe
real futuro de 1,80, una brecha que evidencia \textit{overfitting} o error de
estimaci\'on: la cartera parece excepcional en entrenamiento, pero no replica esa
eficiencia fuera de muestra.

\subsection*{Implicancias y conclusiones}

La primera implicancia es que Markowitz es altamente sensible a la estimaci\'on
de retornos esperados. La matriz de covarianzas importa, pero $\mu$ suele ser
m\'as ruidoso; por eso los perfiles agresivos, que explotan diferencias
peque\~nas entre activos, son los que m\'as se deterioran al incluir pandemia. En
segundo lugar, el Benchmark top-20 equiponderado es una l\'inea base exigente:
diversifica en empresas grandes sin depender de par\'ametros estimados, por lo
que obliga a justificar cualquier superioridad de Markowitz con evidencia fuera
de muestra. En tercer lugar, M\'inima varianza cumple un rol defensivo relevante,
porque no maximiza retorno esperado y depende menos de predicciones individuales.

El rebalanceo semestral debe interpretarse como un \textit{trade-off}.
Actualizar pesos permite incorporar informaci\'on reciente, pero puede amplificar
ruido si las nuevas estimaciones son inestables. Asimismo, la inclusi\'on de
pandemia no significa que el mercado futuro evaluado sea peor, porque el periodo
fuera de muestra se mantiene fijo en ambos escenarios. Lo que cambia es el
conjunto de datos usado para entrenar, por lo que el resultado mide c\'omo los
datos de estr\'es modifican la percepci\'on estimada de riesgo-retorno.

En s\'intesis, el notebook implementa una comparaci\'on metodol\'ogicamente
completa entre Markowitz y benchmark, incorporando escenarios de pandemia,
rebalanceo y validaci\'on por m\'ultiples splits. La principal conclusi\'on es que
no basta con reportar el portafolio de mayor Sharpe en un \'unico periodo: se
debe contrastar con benchmarks simples, evaluar robustez temporal y separar
desempe\~no te\'orico de desempe\~no realizado. El caso \texttt{h7\_f2} queda
justificado por la validaci\'on robusta, mientras que \texttt{h5\_f2} y
\texttt{h6\_f2} deben reportarse como sensibilidad favorable. Para la entrega,
conviene presentar Benchmark y M\'inima varianza como l\'ineas base obligatorias,
usar \texttt{h7\_f2} como escenario principal y evitar afirmar superioridad de
Markowitz sin evidencia fuera de muestra.

\clearpage

% ============================================================
\section*{Black-Litterman}
% ============================================================

\subsection*{Resumen}

Este apartado resume la implementaci\'on del notebook
\texttt{black\_litterman.ipynb}, que extiende el experimento de optimizaci\'on de
portafolios desarrollado previamente con Markowitz. Su objetivo no es reemplazar
la l\'ogica anterior, sino construir una versi\'on Black-Litterman comparable,
usando la misma base de datos, ventanas de entrenamiento y prueba, perfiles de
riesgo e indicadores de evaluaci\'on. De esta forma, las diferencias observadas
se atribuyen principalmente al cambio en la estimaci\'on de retornos esperados:
desde promedios hist\'oricos puros hacia retornos posteriores que combinan
equilibrio de mercado y \textit{views} del inversionista.

El universo central corresponde a 601 activos del mercado F5, seleccionados
desde metadata y series hist\'oricas de retornos totales. Se comparan dos
escenarios: una base sin el periodo pand\'emico y otra que incorpora las
observaciones de 2020--2022. El experimento principal utiliza 7 a\~nos de
calibraci\'on y 2 a\~nos de evaluaci\'on fuera de muestra, configuraci\'on denotada
\texttt{h7\_f2}. Durante el periodo futuro, el portafolio se recalibra
semestralmente cada 126 d\'ias burs\'atiles; en cada rebalanceo se estima la matriz
de covarianzas, se construye el prior Black-Litterman, se aplica una
\textit{view} y se optimizan portafolios con CVXPY.

\begin{table}[H]
\centering
\caption{Escenarios evaluados en Black-Litterman.}
\label{tab:anexo4_bl_escenarios}
\small
\begin{tabular}{@{}p{0.20\textwidth}p{0.33\textwidth}p{0.37\textwidth}@{}}
\toprule
Escenario & Descripci\'on & Prop\'osito \\
\midrule
\texttt{sin\_pandemia} & Base filtrada sin 2020--2022. & Evaluar desempe\~no en una muestra menos dominada por shocks extremos. \\
\texttt{con\_pandemia} & Base que incorpora 2020--2022. & Medir sensibilidad ante un historial con estr\'es, recuperaci\'on y alta volatilidad. \\
\bottomrule
\end{tabular}
\end{table}

\subsection*{Fundamentos metodol\'ogicos}

El modelo Black-Litterman responde a una debilidad central de Markowitz: los
retornos esperados son dif\'iciles de estimar y peque\~nos cambios en ellos pueden
generar variaciones importantes en los pesos \'optimos. Para reducir esa
sensibilidad, Black-Litterman combina dos fuentes de informaci\'on: un prior de
equilibrio de mercado, que representa retornos impl\'icitos en una cartera de
mercado, y \textit{views} del inversionista, que expresan opiniones absolutas o
relativas sobre activos, grupos o factores.

La implementaci\'on calcula el prior como:

\begin{equation}
\pi = \delta \Sigma w_{\text{market}},
\end{equation}

donde $\Sigma$ es la matriz de covarianzas diaria estimada, $w_{\text{market}}$
son pesos por capitalizaci\'on burs\'atil y $\delta$ es la aversi\'on al riesgo
impl\'icita:

\begin{equation}
\delta =
\frac{R_{\text{mercado}} - R_f}{\sigma^2_{\text{mercado}}}.
\end{equation}

Las \textit{views} se representan mediante:

\begin{equation}
P\mu = Q + \varepsilon, \qquad \varepsilon \sim \mathcal{N}(0,\Omega),
\end{equation}

donde $P$ define la combinaci\'on de activos observada, $Q$ el retorno esperado
de esa combinaci\'on y $\Omega$ la incertidumbre asociada. La actualizaci\'on
posterior utilizada es:

\begin{equation}
\mu_{\text{BL}} =
\pi + \tau\Sigma P^\top
(P\tau\Sigma P^\top + \Omega)^{-1}
(Q - P\pi).
\end{equation}

El t\'ermino $Q - P\pi$ representa la sorpresa de la \textit{view} respecto del
prior. Si la opini\'on ya est\'a incorporada en el equilibrio, el ajuste es
peque\~no; si contradice al prior y tiene baja incertidumbre, el posterior se
desplaza con mayor fuerza hacia la \textit{view}.

\noindent\textbf{Construcci\'on y peso de las \textit{views}.}
La \textit{view} de momentum se basa en persistencia de retornos. El c\'odigo
calcula el rendimiento acumulado de cada activo durante los \'ultimos 252 d\'ias,
selecciona los 20 con mayor momentum y los 20 con peor momentum, y construye una
opini\'on relativa:

\begin{equation}
R(\text{top 20 momentum}) - R(\text{bottom 20 momentum}) = Q.
\end{equation}

La matriz $P$ asigna peso positivo igualitario a los ganadores y negativo
igualitario a los rezagados. Esta formulaci\'on es adecuada para Black-Litterman
porque no exige estimar el retorno absoluto de cada activo, sino solo expresar
que un grupo deber\'ia superar a otro.

La \textit{view} de desempleo se basa en sensibilidad c\'iclica sectorial. Si el
desempleo asumido est\'a por debajo de su nivel neutral, la econom\'ia se
interpreta como m\'as fuerte, por lo que sectores c\'iclicos deber\'ian superar a
sectores defensivos. Con los par\'ametros actuales:

\begin{equation}
\text{signal} = 0,05 - 0,04 = 0,01,
\qquad
Q = \frac{0,01}{252} = 0,00003968.
\end{equation}

En consecuencia, $P$ queda long sectores c\'iclicos y short sectores defensivos.
Esta \textit{view} es conceptualmente razonable, pero debe leerse como un
an\'alisis de escenario, ya que el desempleo usado es fijo y no una observaci\'on
actualizada en cada rebalanceo.

El peso de una \textit{view} no se define como porcentaje de portafolio, sino
mediante su precisi\'on estad\'istica relativa frente al prior. En el c\'odigo:

\begin{equation}
\Omega = \frac{P\tau\Sigma P^\top}{\text{confidence}}.
\end{equation}

Si la confianza aumenta, $\Omega$ disminuye y el posterior se mueve con mayor
fuerza hacia $Q$. En esta implementaci\'on, momentum tiene confianza 0,50 y
desempleo 0,35; por tanto, a igualdad de varianza de la \textit{view}, momentum
recibe mayor peso estad\'istico. Sin embargo, el impacto final tambi\'en depende
de $Q$, de $P$, de $\tau$, de $\Sigma$ y de la distancia entre $Q$ y $P\pi$.

\subsection*{Supuestos del experimento}

El experimento utiliza retornos totales diarios, 252 d\'ias burs\'atiles por
a\~no, covarianza con \textit{shrinkage} de 20\%, correcci\'on
\texttt{nearest\_psd}, prior de equilibrio basado en capitalizaci\'on, tasa libre
de riesgo fija de 2\% anual y $\tau = 0,05$. Las carteras finales son
\textit{long-only}, invierten todo el capital, aplican cotas de concentraci\'on y
se rebalancean semestralmente. Las \textit{views} pueden ser long-short, pero
los portafolios finales no permiten ventas cortas. El benchmark top-20
equiponderado no depende de Black-Litterman ni de las \textit{views}.

\begin{table}[H]
\centering
\caption{Supuestos principales del experimento Black-Litterman.}
\label{tab:anexo4_bl_supuestos}
\small
\begin{adjustbox}{max width=\textwidth}
\begin{tabular}{@{}p{0.28\textwidth}p{0.28\textwidth}p{0.34\textwidth}@{}}
\toprule
Supuesto & Implementaci\'on & Implicancia \\
\midrule
Retornos diarios & Series de retornos totales diarios & Se asume frecuencia diaria y 252 d\'ias burs\'atiles. \\
Covarianza estable & \texttt{estimate\_parameters} con 20\% \textit{shrinkage} & Mejora estabilidad, pero no modela cambios de r\'egimen. \\
Matriz PSD & \texttt{nearest\_psd} & Evita problemas num\'ericos en optimizaci\'on cuadr\'atica. \\
Prior de equilibrio & $\pi = \delta\Sigma w_{\text{market}}$ & Supone que capitalizaci\'on contiene informaci\'on agregada. \\
Tasa libre fija & 2\% anual & Simplifica Sharpe y prior. \\
\texttt{tau} & 0,05 & Incertidumbre del prior no calibrada emp\'iricamente. \\
\textit{Long-only} & $w \geq 0$ & Las carteras finales no hacen ventas cortas. \\
Rebalanceo & 126 d\'ias & Actualiza modelo, pero no incorpora costos de transacci\'on. \\
Momentum persistente & Ganadores recientes superan perdedores & Puede ser sensible a reversals. \\
Desempleo fijo & 4\% vs 5\% neutral & Escenario macro, no serie observada din\'amica. \\
\bottomrule
\end{tabular}
\end{adjustbox}
\end{table}

\subsection*{Estructura del notebook}

El notebook importa funciones del m\'odulo
\texttt{markowitz\_cvxpy\_benchmark\_teorico.py}, reutilizando la carga de
datos, estimaci\'on de par\'ametros, perfiles de riesgo, restricciones, KPIs y
optimizadores de Markowitz. Esta decisi\'on permite que la comparaci\'on entre
Markowitz y Black-Litterman sea metodol\'ogicamente consistente.

\begin{table}[H]
\centering
\caption{Par\'ametros globales del experimento Black-Litterman.}
\label{tab:anexo4_bl_parametros_globales}
\small
\begin{adjustbox}{max width=\textwidth}
\begin{tabular}{@{}p{0.34\textwidth}p{0.18\textwidth}p{0.38\textwidth}@{}}
\toprule
Par\'ametro & Valor & Funci\'on \\
\midrule
\texttt{MARKET\_SIZE} & 601 & N\'umero de activos del universo de mercado. \\
\texttt{BENCHMARK\_SIZE} & 20 & N\'umero de activos del benchmark equiponderado. \\
\texttt{HISTORICAL\_YEARS} & 7 & A\~nos de entrenamiento del experimento central. \\
\texttt{FUTURE\_YEARS} & 2 & A\~nos de evaluaci\'on fuera de muestra. \\
\texttt{REBALANCE\_DAYS} & 126 & Frecuencia semestral de rebalanceo. \\
\texttt{DEFAULT\_MAX\_WEIGHT} & 20\% & Cota m\'axima general para portafolios especiales. \\
\texttt{DEFAULT\_RISK\_FREE\_RATE} & 2\% anual & Tasa usada en Sharpe y prior. \\
\texttt{DEFAULT\_SHRINKAGE} & 20\% & Regularizaci\'on de la matriz de covarianzas. \\
\texttt{DEFAULT\_TAU} & 0,05 & Incertidumbre del prior Black-Litterman. \\
\bottomrule
\end{tabular}
\end{adjustbox}
\end{table}

\begin{table}[H]
\centering
\caption{Par\'ametros de las \textit{views}.}
\label{tab:anexo4_bl_parametros_views}
\small
\begin{adjustbox}{max width=\textwidth}
\begin{tabular}{@{}p{0.18\textwidth}p{0.34\textwidth}p{0.16\textwidth}p{0.22\textwidth}@{}}
\toprule
\textit{View} & Par\'ametro & Valor & Interpretaci\'on \\
\midrule
Momentum & \texttt{MOMENTUM\_LOOKBACK} & 252 d\'ias & Ventana reciente para medir momentum. \\
Momentum & \texttt{MOMENTUM\_TOP\_BOTTOM} & 20 & Activos long y short de la \textit{view} relativa. \\
Momentum & \texttt{MOMENTUM\_CONFIDENCE} & 0,50 & Confianza asignada a la \textit{view}. \\
Desempleo & \texttt{UNEMPLOYMENT\_RATE\_ASSUMED} & 4\% & Desempleo asumido. \\
Desempleo & \texttt{UNEMPLOYMENT\_NEUTRAL} & 5\% & Nivel neutral de desempleo. \\
Desempleo & \texttt{MACRO\_UNEMPLOYMENT\_BETA} & 1,0 & Sensibilidad del diferencial macro. \\
Desempleo & \texttt{UNEMPLOYMENT\_CONFIDENCE} & 0,35 & Confianza asignada a la \textit{view} macro. \\
\bottomrule
\end{tabular}
\end{adjustbox}
\end{table}

La clase \texttt{ScenarioDataset} agrupa la llave del escenario, etiqueta
legible, matriz de retornos, metadata y tickers cargados. La funci\'on
\texttt{load\_scenario\_dataset} lee la metadata, selecciona activos de mercado
y concatena retornos con intersecci\'on temporal mediante \texttt{join="inner"},
garantizando una matriz rectangular para optimizaci\'on. La funci\'on
\texttt{market\_cap\_weights} construye pesos de mercado con capitalizaci\'on
burs\'atil y cae a un portafolio equiponderado si las capitalizaciones son
inv\'alidas.

El pipeline central estima par\'ametros con \textit{shrinkage}, calcula el prior
de mercado, construye la \textit{view}, aplica Black-Litterman y resuelve
portafolios para cinco perfiles de riesgo, \texttt{b\_teorico}, M\'inima
varianza, Benchmark e Ideal de mercado. Cuando se ejecutan validaciones m\'as
costosas, se usa una versi\'on reducida que resuelve solo Neutro, M\'inima
varianza y Benchmark. El \textit{smoke test} confirma que $\mu_{\text{BL}}$ y
$\Sigma_{\text{BL}}$ no contienen valores inv\'alidos, que los pesos suman 1 y
que no aparecen posiciones negativas relevantes.

\begin{table}[H]
\centering
\caption{Dimensi\'on de salidas del experimento Black-Litterman.}
\label{tab:anexo4_bl_salidas}
\small
\begin{tabular}{@{}lr@{}}
\toprule
Objeto & Filas \\
\midrule
Resultados centrales rebalanceados & 36 \\
Segmentos centrales & 144 \\
\textit{Views} registradas & 16 \\
Ventanas aleatorias procesadas & 70 \\
Ventanas robustas procesadas & 18 \\
Filas de validaci\'on aleatoria & 420 \\
Filas de validaci\'on robusta & 108 \\
\bottomrule
\end{tabular}
\end{table}

\subsection*{KPIs utilizados}

Los KPIs utilizados son retorno total rebalanceado, retorno anual rebalanceado,
volatilidad diaria y anual, Sharpe rebalanceado, m\'aximo \textit{drawdown},
CVaR diario al 95\%, n\'umero de segmentos, activos con peso positivo y
principales posiciones. El indicador principal para ordenar resultados es el
Sharpe rebalanceado, porque resume retorno y riesgo fuera de muestra tras
concatenar todos los rebalanceos.

\subsection*{Resultados del experimento pandemia}

Los mejores portafolios Black-Litterman del experimento central se muestran en
la Tabla~\ref{tab:anexo4_bl_mejores}. El mejor resultado global corresponde a
la \textit{view} de desempleo sin pandemia en el portafolio Neutro, con Sharpe
2,44. Momentum tambi\'en funciona bien sin pandemia, aunque pierde fuerza
relativa cuando se incluye el periodo pand\'emico.

\begin{table}[H]
\centering
\caption{Mejores portafolios Black-Litterman del experimento central.}
\label{tab:anexo4_bl_mejores}
\scriptsize
\begin{adjustbox}{max width=\textwidth}
\begin{tabular}{@{}lllrrrr@{}}
\toprule
Escenario & \textit{View} & Portafolio l\'ider & Ret. total & Vol. anual & Max drawdown & Sharpe \\
\midrule
Sin pandemia & Desempleo & Neutro & 63,1\% & 10,5\% & $-$10,4\% & 2,44 \\
Sin pandemia & Momentum & Muy conservador & 64,0\% & 10,9\% & $-$10,1\% & 2,38 \\
Con pandemia & Desempleo & Arriesgado & 71,9\% & 12,4\% & $-$13,3\% & 2,36 \\
Con pandemia & Momentum & Conservador & 56,4\% & 11,2\% & $-$10,3\% & 2,06 \\
\bottomrule
\end{tabular}
\end{adjustbox}
\end{table}

\begin{table}[H]
\centering
\caption{Comparaci\'on final entre equiponderado, Markowitz y Black-Litterman.}
\label{tab:anexo4_bl_comparacion_final}
\scriptsize
\begin{adjustbox}{max width=\textwidth}
\begin{tabular}{@{}llllrrrrr@{}}
\toprule
Escenario & Modelo & \textit{View} & Portafolio & Ret. total & Vol. anual & Drawdown & CVaR 95\% & Sharpe \\
\midrule
Con pandemia & Equiponderado & No aplica & Benchmark & 101,4\% & 19,5\% & $-$21,9\% & 2,79\% & 2,05 \\
Con pandemia & Markowitz & Hist\'orico & Muy conservador & 54,5\% & 11,0\% & $-$10,2\% & 1,50\% & 2,02 \\
Con pandemia & Black-Litterman & Desempleo & Arriesgado & 71,9\% & 12,4\% & $-$13,3\% & 1,71\% & 2,36 \\
Con pandemia & Black-Litterman & Momentum & Conservador & 56,4\% & 11,2\% & $-$10,3\% & 1,47\% & 2,06 \\
Sin pandemia & Equiponderado & No aplica & Benchmark & 101,4\% & 19,5\% & $-$21,9\% & 2,79\% & 2,05 \\
Sin pandemia & Markowitz & Hist\'orico & Ideal de mercado & 76,5\% & 12,3\% & $-$12,6\% & 1,67\% & 2,51 \\
Sin pandemia & Black-Litterman & Desempleo & Neutro & 63,1\% & 10,5\% & $-$10,4\% & 1,41\% & 2,44 \\
Sin pandemia & Black-Litterman & Momentum & Muy conservador & 64,0\% & 10,9\% & $-$10,1\% & 1,45\% & 2,38 \\
\bottomrule
\end{tabular}
\end{adjustbox}
\end{table}

La lectura principal es que Black-Litterman mejora claramente a Markowitz en el
escenario con pandemia, especialmente con la \textit{view} de desempleo. En el
escenario sin pandemia, Markowitz alcanza el mayor Sharpe final con el
portafolio Ideal de mercado, aunque Black-Litterman obtiene resultados cercanos
con menor volatilidad y \textit{drawdowns} controlados.

\begin{table}[H]
\centering
\caption{Mayor mejora de Black-Litterman frente a Markowitz.}
\label{tab:anexo4_bl_delta_markowitz}
\small
\begin{tabular}{@{}lllrrr@{}}
\toprule
Escenario & \textit{View} & Portafolio & Sharpe BL & Sharpe Markowitz & Delta Sharpe \\
\midrule
Con pandemia & Desempleo & Arriesgado & 2,36 & 1,17 & +1,19 \\
\bottomrule
\end{tabular}
\end{table}

Tambi\'en destacan las mejoras de desempleo/Muy arriesgado con delta Sharpe
+0,87, desempleo/\texttt{b\_teorico} con +0,64 y momentum/Muy arriesgado con
+0,59. Esto sugiere que Black-Litterman aporta mayor valor cuando el historial
contiene episodios extremos, ya que el prior de equilibrio y las \textit{views}
reducen la dependencia de medias hist\'oricas inestables.

Al comparar el mismo portafolio y \textit{view} entre bases con y sin pandemia,
varios deltas de Sharpe son negativos. Momentum/Muy conservador cae $-0,37$,
desempleo/Neutro cae $-0,31$, desempleo/Muy conservador cae $-0,30$ y
desempleo/Conservador cae $-0,29$. Esto indica que incluir pandemia tiende a
deteriorar la estabilidad relativa de varias configuraciones, pero tambi\'en
muestra que el escenario con pandemia es una prueba m\'as exigente donde ciertas
\textit{views}, especialmente desempleo, siguen mejorando a Markowitz.

\subsection*{Rebalanceo}

El rebalanceo Black-Litterman mantiene la frecuencia semestral de 126 d\'ias
burs\'atiles. En cada fecha se reestiman covarianzas, prior de equilibrio,
\textit{views}, retornos posteriores y pesos finales. Esta decisi\'on permite que
la se\~nal incorpore informaci\'on disponible en cada segmento, pero mantiene la
comparabilidad con Markowitz porque usa la misma ventana futura, las mismas
restricciones y los mismos KPIs. Como en Markowitz, el rebalanceo no incorpora
costos de transacci\'on, por lo que sus resultados deben interpretarse como
desempe\~no te\'orico fuera de muestra.

\subsection*{Splits aleatorios}

En los 70 splits aleatorios, los mejores resultados para el portafolio Neutro
favorecen horizontes con 6 a\~nos de entrenamiento y 2 de prueba. Esto es
consistente con la necesidad de contar con suficiente historia para estabilizar
covarianzas y prior, sin usar una ventana tan larga que diluya cambios
estructurales.

\begin{table}[H]
\centering
\caption{Validaci\'on con 70 splits aleatorios para el portafolio Neutro.}
\label{tab:anexo4_bl_splits}
\small
\begin{tabular}{@{}lrrrrr@{}}
\toprule
\textit{View} & Horizonte & Sharpe prom. & Ret. total prom. & Vol. anual prom. & Ventanas \\
\midrule
Desempleo & \texttt{h6\_f2} & 2,48 & 61,5\% & 10,4\% & 3 \\
Momentum & \texttt{h6\_f2} & 2,31 & 68,6\% & 13,1\% & 3 \\
\bottomrule
\end{tabular}
\end{table}

\subsection*{Evaluaci\'on robusta}

\begin{table}[H]
\centering
\caption{Validaci\'on robusta con 18 ventanas para el portafolio Neutro.}
\label{tab:anexo4_bl_validacion_robusta}
\small
\begin{tabular}{@{}lrrrrr@{}}
\toprule
\textit{View} & Mejor horizonte & Sharpe prom. & Ret. total prom. & Vol. anual prom. & Drawdown prom. \\
\midrule
Desempleo & \texttt{h7\_f2} & 2,35 & 59,2\% & 10,6\% & $-$9,7\% \\
Momentum & \texttt{h7\_f2} & 2,23 & 68,5\% & 13,6\% & $-$12,8\% \\
\bottomrule
\end{tabular}
\end{table}

En validaci\'on robusta, desempleo domina en Sharpe para el portafolio Neutro
porque logra menor volatilidad y menor \textit{drawdown}. Momentum obtiene
mayor retorno total promedio, pero a costa de mayor riesgo.

\subsection*{Implicancias y conclusiones}

Los resultados muestran que Black-Litterman cumple su objetivo principal:
suavizar la dependencia de medias hist\'oricas puras. Esto se observa con mayor
claridad en el escenario con pandemia, donde Markowitz queda m\'as castigado y
Black-Litterman mejora el Sharpe en varios portafolios. El tipo de \textit{view}
es determinante. Desempleo genera portafolios m\'as defensivos y mejor Sharpe en
configuraciones centrales; aunque conceptualmente es una \textit{view}
c\'iclica, su menor $Q$ y menor confianza producen ajustes m\'as moderados.
Momentum, en cambio, entrega una se\~nal m\'as intensa y puede capturar retornos
altos, pero tambi\'en expone a reversals y mayor volatilidad.

El benchmark equiponderado top-20 obtiene retornos totales elevados, aunque con
volatilidad y \textit{drawdown} superiores. Su Sharpe es competitivo, pero no
domina siempre, lo que muestra que un benchmark simple puede ser dif\'icil de
superar en retorno bruto mientras los modelos optimizados pueden mejorar la
relaci\'on retorno-riesgo. Las diferencias entre \texttt{sin\_pandemia} y
\texttt{con\_pandemia} deben entenderse como una prueba de robustez: incluir
shocks hist\'oricos cambia la frontera eficiente y las estimaciones de riesgo, no
necesariamente el mercado futuro evaluado.

Para mejorar la implementaci\'on, la \textit{view} de desempleo deber\'ia pasar de
supuesto fijo a se\~nal observable o pronosticada, usando por ejemplo la tasa
efectiva disponible al momento del rebalanceo, sorpresas respecto de
expectativas o variaciones mensuales o trimestrales. La \textit{view} de
momentum podr\'ia refinarse con momentum ajustado por volatilidad, exclusi\'on del
mes m\'as reciente para reducir reversals o combinaci\'on de ventanas de 6 y
12 meses. Tambi\'en ser\'ia valioso incorporar \textit{views} de valor, calidad,
tama\~no, tasas de inter\'es, inflaci\'on o revisiones de utilidades.

Otra mejora importante es calibrar $\Omega$ con evidencia emp\'irica en lugar de
fijar manualmente la confianza. Esta calibraci\'on podr\'ia usar error hist\'orico
de predicci\'on, frecuencia de acierto del signo, correlaci\'on entre se\~nal y
retorno futuro, varianza residual de regresiones, metodolog\'ia de
desplazamiento deseado en pesos o validaci\'on cruzada temporal. Adem\'as,
conviene analizar sensibilidad a $\tau$, \textit{shrinkage}, frecuencia de
rebalanceo, \texttt{max\_weight} y tasa libre de riesgo, e incorporar costos de
transacci\'on, restricciones sectoriales, cotas de liquidez, control de
\textit{tracking error}, penalizaci\'on por rotaci\'on y estimadores robustos de
covarianza.

En s\'intesis, la implementaci\'on construye un experimento Black-Litterman
completo y comparable con Markowitz y benchmark. Su principal aporte es
reemplazar retornos esperados hist\'oricos por retornos posteriores que combinan
equilibrio de mercado y \textit{views}. Los resultados indican que
Black-Litterman es especialmente valioso en escenarios exigentes, como el
entrenamiento con pandemia, donde la \textit{view} de desempleo mejora
fuertemente el Sharpe frente a Markowitz. La conclusi\'on central no es que una
\textit{view} sea universalmente superior, sino que la calidad, calibraci\'on y
confianza de las \textit{views} determinan el valor pr\'actico del modelo.

\clearpage

% ============================================================
\section*{Simulaci\'on Monte Carlo}
% ============================================================

\subsection*{Resumen}

Como parte de la documentaci\'on t\'ecnica del proyecto FinPUC, se elabor\'o un
informe consolidado que integra los resultados de los tres notebooks de
implementaci\'on: Markowitz vs Benchmark, Black-Litterman y Simulaci\'on Monte
Carlo. El prop\'osito del documento es proveer trazabilidad completa entre los
datos, c\'alculos, t\'ecnicas de optimizaci\'on, comparaciones y conclusiones
finales del proyecto.

El informe fue generado el 09 de mayo de 2026 y opera sobre el universo F5 de
601 activos, benchmark top-20, rebalanceo semestral de 126 d\'ias h\'abiles y
simulaci\'on Monte Carlo a 260 semanas. En esta etapa, el notebook
\texttt{p4\_montecarlo\_recomendador\_portafolios.ipynb} transforma los KPIs de
Markowitz, Black-Litterman y Equiponderado en una evaluaci\'on de cliente y
empresa. Para cada combinaci\'on de escenario, t\'ecnica, \textit{view} y perfil
simula trayectorias de riqueza, probabilidad de retiro, aceptaci\'on de
recomendaciones, utilidad de la empresa y un score final de recomendaci\'on.

\subsection*{Fundamentos metodol\'ogicos}

La Simulaci\'on Monte Carlo usa como insumos el retorno anual base y la
volatilidad anual de cada estrategia. Estos par\'ametros provienen de los
experimentos previos y se transforman a frecuencia semanal para simular
trayectorias de riqueza de un cliente con capital inicial $C_0 = 1\,000$ USD. El
retiro depende de la p\'erdida relativa frente a la tolerancia del perfil,
mientras que la aceptaci\'on depende del atractivo riesgo-retorno de la
recomendaci\'on.

El informe documenta una metodolog\'ia de c\'alculo com\'un para los modelos de
entrada: estimaci\'on de par\'ametros ($\mu$, $\Sigma$ con \textit{shrinkage}
$\lambda = 0.20$, \texttt{nearest\_psd}), anualizaci\'on y frontera eficiente.
La simulaci\'on no reoptimiza portafolios, sino que traduce esos KPIs financieros
a variables de experiencia de cliente y de resultado para la empresa.

\subsection*{Supuestos del experimento}

Se simulan 5\,000 trayectorias por combinaci\'on durante 260 semanas, equivalentes
a cinco a\~nos de evaluaci\'on semanal. Cada combinaci\'on se eval\'ua en tres
escenarios proyectados: desfavorable, neutro y favorable. Los retornos semanales
se aproximan con distribuci\'on normal, pese a que el an\'alisis distribucional
muestra colas pesadas. La comisi\'on inicial es $k = 1\%$ sobre el capital
inicial y, en recomendaciones aceptadas, se cobra sobre una fracci\'on de
rotaci\'on de 5\% de la riqueza vigente.

El dise\~no opera sobre el mismo universo F5 de 601 activos usado por los
experimentos de optimizaci\'on, benchmark top-20 y KPIs rebalanceados. Por lo
tanto, la comparaci\'on final no cambia el universo de acciones, sino la manera
en que se transforma el desempe\~no hist\'orico de cada t\'ecnica en una
distribuci\'on de resultados de negocio.

\subsection*{Estructura del notebook}

El documento metodol\'ogico consolidado se organiza en ocho secciones:

\begin{enumerate}
    \item \textbf{Resumen ejecutivo:} comparaci\'on final de los tres modelos por
    escenario, con tabla de KPIs y top 12 del ranking.
    \item \textbf{Datos, universo F5 y dise\~no experimental:} descripci\'on del
    filtrado (1\,406 a 601 tickers), split \texttt{h7\_f2} y escenarios con/sin
    pandemia.
    \item \textbf{Metodolog\'ia de c\'alculo com\'un:} estimaci\'on de par\'ametros
    ($\mu$, $\Sigma$ con \textit{shrinkage} $\lambda = 0.20$,
    \texttt{nearest\_psd}), anualizaci\'on y frontera eficiente.
    \item \textbf{Markowitz vs benchmark:} portafolios por perfil, impacto de la
    pandemia en Sharpe, validaci\'on robusta por horizonte.
    \item \textbf{Black-Litterman y matriz de views:} configuraci\'on de las
    cuatro views (momentum, momentum top-20 6M, momentum top-20/bottom-20 1Y,
    desempleo), resultados por view y escenario.
    \item \textbf{Simulaci\'on Monte Carlo y recomendador:} supuestos operativos,
    promedios por t\'ecnica, ranking completo del recomendador.
    \item \textbf{Conclusi\'on integrada:} hallazgos principales y recomendaci\'on
    final.
    \item \textbf{\'Indice de anexos y tablas de auditor\'ia:} inventario
    completo de CSVs generados.
\end{enumerate}

El notebook define funciones de retiro y aceptaci\'on, construye insumos por
t\'ecnica y perfil desde los CSV de Markowitz y Black-Litterman, ejecuta la
Simulaci\'on Monte Carlo, resume KPIs por escenario, genera gr\'aficos, calcula el
ranking recomendador cliente-empresa y exporta conclusiones.

Las salidas principales son:
\begin{flushleft}
\footnotesize
\path|p4_model_inputs_kpis.csv|\\
\path|p4_montecarlo_summary.csv|\\
\path|p4_neutral_scenario_summary.csv|\\
\path|p4_recommendation_ranking.csv|\\
\path|p4_montecarlo_paths_quantiles.csv|
\end{flushleft}

El informe documenta un total de 21 archivos CSV generados por los tres
notebooks, organizados por metodolog\'ia. La
Tabla~\ref{tab:anexo4_csv_index} presenta el inventario completo.

\begin{table}[H]
\centering
\caption{\'Indice de CSVs generados por el proyecto.}
\label{tab:anexo4_csv_index}
\footnotesize
\begin{adjustbox}{max width=\textwidth}
\begin{tabular}{@{}p{0.45\textwidth}lrrl@{}}
\toprule
Archivo & Metodolog\'ia & Filas & Cols & Prop\'osito \\
\midrule
rebalance\_summary\_h7\_f2\_601activos & Markowitz & 18 & 12 & KPIs rebalanceados \\
comparison\_sin\_vs\_con\_pandemia\_h7\_f2 & Markowitz & 9 & 21 & Comparaci\'on escenarios \\
validacion\_robusta\_ventanas\_moviles & Markowitz & 6 & 4 & Sharpe por horizonte \\
frontier\_segments\_h7\_f2\_601activos & Markowitz & 238 & 24 & Frontera eficiente \\
heatmap\_sharpe\_train\_test & Markowitz & 7 & 8 & Heatmap splits \\
\midrule
bl\_rebalance\_summary\_h7\_f2\_601activos & Black-Litterman & 72 & 14 & KPIs BL rebalanceados \\
bl\_views\_h7\_f2\_601activos & Black-Litterman & 32 & 36 & Configuraci\'on views \\
comparison\_three\_models\_kpis\_h7\_f2 & Black-Litterman & 12 & 13 & Comparaci\'on 3 modelos \\
comparison\_bl\_vs\_markowitz\_h7\_f2 & Black-Litterman & 72 & 19 & Delta BL vs Markowitz \\
bl\_robust\_validation\_summary & Black-Litterman & 36 & 10 & Validaci\'on robusta BL \\
\midrule
p4\_montecarlo\_summary & Simulaci\'on Monte Carlo & 180 & 29 & Resumen completo MC \\
p4\_neutral\_scenario\_summary & Simulaci\'on Monte Carlo & 60 & 29 & Escenario neutro \\
p4\_recommendation\_ranking & Simulaci\'on Monte Carlo & 60 & 30 & Ranking recomendador \\
p4\_model\_inputs\_kpis & Simulaci\'on Monte Carlo & 60 & 13 & Inputs desde modelos \\
p4\_montecarlo\_paths\_quantiles & Simulaci\'on Monte Carlo & 1\,980 & 13 & Bandas P5--P95 \\
\bottomrule
\end{tabular}
\end{adjustbox}
\end{table}

\subsection*{KPIs utilizados}

Los KPIs de cliente son riqueza terminal media, percentiles 5/50/95, ganancia
esperada, probabilidad de ganancia, probabilidad de p\'erdida mayor a la
tolerancia, tasa de retiro y semana promedio de retiro. Los KPIs de empresa son
ingreso promedio, ingreso mediano, ingreso percentil 95, ingreso como porcentaje
del capital inicial y tasa de aceptaci\'on de recomendaciones. El ranking final
usa \texttt{p4\_score}, que combina valor para el cliente y valor para la
empresa.

En el informe integrado, estos KPIs se conectan con retorno total, retorno
anual, volatilidad anual, \textit{drawdown}, CVaR y Sharpe de los modelos de
entrada. Por esta raz\'on, una t\'ecnica con mayor retorno hist\'orico puede perder
en score si eleva demasiado retiros o dispersi\'on de resultados.

\subsection*{Resultados del experimento pandemia}

La Tabla~\ref{tab:anexo4_comparacion_modelos} presenta la comparaci\'on final de
modelos por escenario, extra\'ida del informe de metodolog\'ias.

\begin{table}[H]
\centering
\caption{Comparaci\'on final de modelos por escenario (mejor portafolio por modelo).}
\label{tab:anexo4_comparacion_modelos}
\small
\begin{adjustbox}{max width=\textwidth}
\begin{tabular}{@{}lllrrrrr@{}}
\toprule
Escenario & Modelo & Portafolio & Ret. total & Ret. anual & Vol. anual & Drawdown & Sharpe \\
\midrule
Sin pandemia & Markowitz & Ideal de mercado & 76,5\% & 32,8\% & 12,3\% & $-$12,6\% & 2,51 \\
Sin pandemia & BL mom. top20 6m & Neutro & 72,3\% & 31,2\% & 11,9\% & $-$12,9\% & 2,46 \\
Sin pandemia & Equiponderado & Benchmark & 101,4\% & 41,9\% & 19,5\% & $-$21,9\% & 2,05 \\
\midrule
Con pandemia & BL mom. top20 6m & Neutro & 69,5\% & 30,2\% & 11,9\% & $-$11,7\% & 2,38 \\
Con pandemia & BL desempleo & Arriesgado & 71,9\% & 31,1\% & 12,4\% & $-$13,3\% & 2,36 \\
Con pandemia & Equiponderado & Benchmark & 101,4\% & 41,9\% & 19,5\% & $-$21,9\% & 2,05 \\
Con pandemia & Markowitz & Muy conservador & 54,5\% & 24,3\% & 11,0\% & $-$10,2\% & 2,02 \\
\bottomrule
\end{tabular}
\end{adjustbox}
\end{table}

La Tabla~\ref{tab:anexo4_ranking} presenta el top 10 del ranking en escenario
neutro.

\begin{table}[H]
\centering
\caption{Top 10 ranking de Simulaci\'on Monte Carlo en escenario neutro.}
\label{tab:anexo4_ranking}
\footnotesize
\begin{adjustbox}{max width=\textwidth}
\begin{tabular}{@{}clllrrrr@{}}
\toprule
Rank & Escenario & T\'ecnica & Perfil & Riqueza & Retiro & Utilidad & Score \\
\midrule
1 & Con pandemia & BL Mom. Top20 6M & Muy arriesgado & USD 8\,974 & 3,1\% & USD 474 & 9\,417 \\
2 & Sin pandemia & BL Mom. Top20 6M & Muy arriesgado & USD 7\,753 & 2,0\% & USD 419 & 8\,151 \\
3 & Con pandemia & BL Momentum & Muy arriesgado & USD 6\,436 & 14,6\% & USD 335 & 6\,625 \\
4 & Sin pandemia & BL Desempleo & Muy arriesgado & USD 5\,434 & 0,5\% & USD 221 & 5\,650 \\
5 & Sin pandemia & Markowitz & Muy arriesgado & USD 5\,349 & 2,1\% & USD 215 & 5\,543 \\
6 & Sin pandemia & BL Momentum & Arriesgado & USD 5\,304 & 10,9\% & USD 320 & 5\,515 \\
7 & Con pandemia & Equiponderado & Muy arriesgado & USD 5\,280 & 0,0\% & USD 204 & 5\,485 \\
8 & Con pandemia & BL Mom. Top20 6M & Arriesgado & USD 5\,137 & 0,5\% & USD 312 & 5\,445 \\
9 & Sin pandemia & Equiponderado & Muy arriesgado & USD 5\,222 & 0,0\% & USD 203 & 5\,425 \\
10 & Sin pandemia & Equiponderado & Arriesgado & USD 5\,084 & 0,2\% & USD 308 & 5\,390 \\
\bottomrule
\end{tabular}
\end{adjustbox}
\end{table}

La mejor combinaci\'on es Black-Litterman Momentum Top20 6M con datos de
pandemia y perfil Muy arriesgado. Su ventaja se explica por alta riqueza
terminal, tasa de retiro baja para un perfil agresivo y mayor utilidad para la
empresa.

\subsection*{Rebalanceo}

La Simulaci\'on Monte Carlo no recalcula pesos de portafolio. Recibe KPIs que ya
incorporan el rebalanceo semestral de Markowitz y Black-Litterman. Por eso, el
rebalanceo entra de forma indirecta: afecta retornos, volatilidad, drawdown y
CVaR base, que luego determinan riqueza simulada, retiro y utilidad. Esta
estructura evita mezclar dos capas de decisi\'on: la optimizaci\'on de pesos queda
en los modelos financieros y la simulaci\'on se concentra en la evaluaci\'on
cliente-empresa.

\subsection*{Splits aleatorios}

La Simulaci\'on Monte Carlo no ejecuta splits aleatorios propios. Su robustez
temporal depende de los insumos previamente validados en Markowitz y
Black-Litterman. Esta decisi\'on evita duplicar experimentos y mantiene el foco
en la pregunta de negocio: cu\'al combinaci\'on de t\'ecnica, escenario y perfil
genera mejor equilibrio cliente-empresa. En consecuencia, la lectura de sus
resultados debe considerar la evidencia de splits aleatorios y validaci\'on
robusta reportada para los modelos de entrada.

\subsection*{Evaluaci\'on robusta}

\begin{table}[H]
\centering
\caption{Ganadores por escenario proyectado en Simulaci\'on Monte Carlo.}
\label{tab:anexo4_mc_robusta}
\small
\begin{adjustbox}{max width=\textwidth}
\begin{tabular}{llllrrr}
\toprule
Escenario & T\'ecnica & View & Perfil & Riqueza & Retiro & Utilidad \\
\midrule
Desfavorable & BL Momentum Top20 6M & momentum\_top20\_6m & Muy arriesgado & USD 5\,009 & 8,8\% & USD 327 \\
Neutro & BL Momentum Top20 6M & momentum\_top20\_6m & Muy arriesgado & USD 8\,974 & 3,1\% & USD 474 \\
Favorable & BL Momentum Top20 6M & momentum\_top20\_6m & Muy arriesgado & USD 15\,451 & 1,6\% & USD 690 \\
\bottomrule
\end{tabular}
\end{adjustbox}
\end{table}

La misma t\'ecnica domina en los tres escenarios proyectados, lo que refuerza su
consistencia dentro del marco de supuestos del notebook. Esta conclusi\'on debe
leerse con cautela porque los retornos simulados son normales y pueden
subestimar eventos extremos. Para implementaci\'on comercial, conviene mostrar al
usuario tres escenarios y no solo el escenario neutro: el riesgo de retiro cambia
fuertemente cuando el camino simulado atraviesa p\'erdidas tempranas.

\subsection*{Implicancias y conclusiones}

El informe de metodolog\'ias concluye con las siguientes observaciones:

\begin{itemize}
    \item El benchmark top-20 equiponderado es una referencia fuerte por retorno
    total, pero su volatilidad y drawdown son mayores que las carteras
    Black-Litterman ajustadas por riesgo.
    \item Markowitz aporta una base estable y trazable, especialmente para
    perfiles conservadores, aunque depende fuertemente de retornos hist\'oricos.
    \item Black-Litterman Desempleo ofrece un ajuste macro competitivo y en
    algunos escenarios lidera frente a Markowitz por delta Sharpe.
    \item Black-Litterman Momentum Top20 6M es la t\'ecnica m\'as fuerte en la
    Simulaci\'on Monte Carlo y la mejor variante market-cap momentum por Sharpe
    en ambos escenarios de datos.
    \item La recomendaci\'on final depende del objetivo: si se prioriza riqueza
    esperada y score, gana Black-Litterman Momentum Top20 6M; si se prioriza
    m\'inimo retiro, se deben revisar perfiles menos agresivos o
    Black-Litterman Desempleo.
\end{itemize}

La Simulaci\'on Monte Carlo conecta desempe\~no financiero con variables de
negocio. El resultado principal es que Black-Litterman Momentum Top20 6M domina
el ranking al combinar alta riqueza esperada, baja tasa de retiro y mayor
utilidad para la empresa. La principal limitaci\'on es el supuesto normal: para
la siguiente etapa conviene probar t-Student, Johnson SU u otra distribuci\'on
con colas pesadas, adem\'as de sensibilizar costos de transacci\'on y reglas de
retiro basadas en \textit{drawdown}.

El informe final de metodolog\'ias de portafolios se encuentra disponible como
archivo PDF en el repositorio del proyecto:
\url{https://github.com/Psacevedo/Capstone.git}. El documento incluye
16 figuras, 18 tablas y el \'indice completo de 21 CSVs de auditor\'ia.

% ============================================================
\section*{Conclusi\'on integrada}
% ============================================================

Los tres notebooks forman una cadena metodol\'ogica. Markowitz define la base y
el benchmark; Black-Litterman reduce dependencia de medias hist\'oricas mediante
prior y \textit{views}; Simulaci\'on Monte Carlo traduce los KPIs en una decisi\'on
cliente-empresa. La recomendaci\'on actual es usar Black-Litterman Momentum Top20
6M como alternativa principal, mantener el benchmark como referencia exigente y
usar Markowitz como base de comparaci\'on y diagn\'ostico de sensibilidad.
